% Generated. Do not edit.
% Projection scriptural-date-location/evidence, version 1: the Scriptural Date and
% Location page, generated whole from research/chronology.toml, the Scripture
% chronology corpus and the source library by tools/tpt document-types refresh.
% document: liturgy/roman-rite/1962/propers/temporal/45-fifth-after-pentecost
\section*{Appendix: Evidence for the Scriptural Dates}
\label{sec:date-location-evidence}
{\small For each passage, as the Scriptural Date and Location table names it: what is dated, how the corpus weighs each answer, the source's own words, and the sources cited. A subject the corpus does not date is reported as such and given no date.\par}
{\small
\par\smallskip\noindent\textit{Offertory: Ps. 15:7--8 (Vg.; Heb. 16)}\par\nobreak
\begin{itemize}[leftmargin=1.2em,itemsep=1pt,topsep=1pt,parsep=0pt]
\item \textbf{Historical setting}. No date: the chronology corpus has not yet dated this (research pending).
\item \textbf{Composition} (preferred). Psalms: Before c.~165 B.C. (``before the Maccabean period, around 165 B.C.; no individual psalm can be dated securely''; Catholic critical chronology; New American Bible, Revised Edition --- USCCB web state (2026), Psalms introduction).
\item \textbf{Prophetic referent} (preferred). The Resurrection of Our Lord: Three days after the Crucifixion (The Catholic Encyclopedia, Volume 3 (New York, 1908)).
\end{itemize}
\par\smallskip\noindent\textit{Alleluia: Ps. 20:1 (Vg.; Heb. 21)}\par\nobreak
\begin{itemize}[leftmargin=1.2em,itemsep=1pt,topsep=1pt,parsep=0pt]
\item \textbf{Historical setting}. As under Offertory: Ps. 15:7--8 (Vg.; Heb. 16).
\item \textbf{Composition}. As under Offertory: Ps. 15:7--8 (Vg.; Heb. 16).
\end{itemize}
\par\smallskip\noindent\textit{Introit: Ps. 26:7, 9 (Vg.; Heb. 27), Ps. 26:1 (Vg.; Heb. 27)}\par\nobreak
\begin{itemize}[leftmargin=1.2em,itemsep=1pt,topsep=1pt,parsep=0pt]
\item \textbf{Historical setting}. As under Offertory: Ps. 15:7--8 (Vg.; Heb. 16).
\item \textbf{Composition}. As under Offertory: Ps. 15:7--8 (Vg.; Heb. 16).
\end{itemize}
\par\smallskip\noindent\textit{Communion: Ps. 26:4 (Vg.; Heb. 27)}\par\nobreak
\begin{itemize}[leftmargin=1.2em,itemsep=1pt,topsep=1pt,parsep=0pt]
\item \textbf{Historical setting}. As under Offertory: Ps. 15:7--8 (Vg.; Heb. 16).
\item \textbf{Composition}. As under Offertory: Ps. 15:7--8 (Vg.; Heb. 16).
\end{itemize}
\par\smallskip\noindent\textit{Gradual: Ps. 83:10 (Vg.; Heb. 84), Ps. 83:9 (Vg.; Heb. 84)}\par\nobreak
\begin{itemize}[leftmargin=1.2em,itemsep=1pt,topsep=1pt,parsep=0pt]
\item \textbf{Historical setting}. As under Offertory: Ps. 15:7--8 (Vg.; Heb. 16).
\item \textbf{Composition} (preferred). The Psalms of the Sons of Korah: Between the days of Isaias and the return from exile (The Catholic Encyclopedia, Volume 12 (1911)).
\end{itemize}
\par\smallskip\noindent\textit{Gospel: Matt. 5:20--24}\par\nobreak
\begin{itemize}[leftmargin=1.2em,itemsep=1pt,topsep=1pt,parsep=0pt]
\item \textbf{Narrated event}. No date: the chronology corpus has not yet dated this (research pending).
\item \textbf{Composition} (disputed). The Gospel of St Matthew: c.~A.D.~38--45 (``about A.D. 38-45''; The Catholic Encyclopedia, Volume 10 (1911); The Catholic Encyclopedia, Volume 14 (New York, 1912)); c.~A.D.~40--42 (``about the year 40-42''; The Catholic Encyclopedia, Volume 10 (1911)); A.D.~40--45 (``the years 40-45''; The Catholic Encyclopedia, Volume 10 (1911)); c.~A.D.~60--68 (``about the year 60-68''; The Catholic Encyclopedia, Volume 10 (1911)); c.~A.D.~64--67 (``about the years 64-67''; The Catholic Encyclopedia, Volume 10 (1911)); c.~A.D.~50 (``about the year 50''; The Catholic Encyclopedia, Volume 14 (New York, 1912)). Place: Semitic original associated with Palestine (The Catholic Encyclopedia, Volume 10 (1911)).
\end{itemize}
\par\smallskip\noindent\textit{Epistle: 1 Pet. 3:8--15}\par\nobreak
\begin{itemize}[leftmargin=1.2em,itemsep=1pt,topsep=1pt,parsep=0pt]
\item \textbf{Composition} (preferred). The First Epistle of St Peter: A.D.~63--64 (``about the end of the year 63 or the beginning of 64''; The Catholic Encyclopedia, Volume 11 (New York, 1911)).
\end{itemize}
\par}
